\documentclass{article}
\usepackage[T1]{fontenc}
\usepackage{lmodern}
\usepackage{graphicx}
\usepackage{amsmath,amssymb}
\usepackage[a4paper,margin=2cm]{geometry}
\usepackage[absolute,overlay]{textpos}
\setlength{\TPHorizModule}{1mm}
\setlength{\TPVertModule}{1mm}
\usepackage[indonesian]{babel}
\usepackage{hyperref}
\setlength{\emergencystretch}{2em}
% unit: Halaman 22

\begin{document}

% === Halaman 22 (llm) ===

\begin{itemize}
    \item Tahun 1977, 3 orang penemu $RSA$ membuat sayembara untuk memecahkan cipherteks dengan menggunakan RSA di majalah \textit{Scientific American}.
    \item Hadiahnya: \$100
    \item Tahun 1994, kelompok yang bekerja dengan kolaborasi internet berhasil memecahkan cipherteks hanya dalam waktu 8 bulan.
\end{itemize}

\end{document}
